\documentclass[11pt,a4paper]{article}

% Compile with pdfLaTeX on Overleaf. The project is self-contained.
\usepackage[T1]{fontenc}
\usepackage[utf8]{inputenc}
\usepackage{lmodern}
\usepackage[margin=1in]{geometry}
\usepackage{amsmath,amssymb,amsthm,mathtools}
\usepackage{booktabs,tabularx,array}
\usepackage{microtype}
\usepackage{needspace}
\usepackage[round,authoryear]{natbib}
\usepackage[hidelinks]{hyperref}

\newcommand{\J}{\mathcal{J}}
\newcommand{\I}{\mathcal{I}}
\newcommand{\Eset}{\mathcal{E}}
\newcommand{\A}{\mathcal{A}}
\newcommand{\Z}{\mathcal{Z}}
\newcommand{\Om}{\Omega}
\newcommand{\EE}{\mathbb{E}}
\newcommand{\PP}{\mathbb{P}}
\newcommand{\ind}{\mathbf{1}}
\DeclareMathOperator{\Var}{Var}
\DeclareMathOperator{\Cov}{Cov}
\numberwithin{equation}{section}
\newtheorem{proposition}{Proposition}

\title{A Probing-Enhanced Stochastic Programming Model for\\
Intra-Zone Homogeneity Delineation and Nutrient Application}
\author{Working draft -- author information to be added}
\date{}

\begin{document}
\maketitle

\begin{abstract}
We propose a probing-enhanced stochastic extension of the Intra-Zone
Homogeneity Delineation Problem for cost-efficient site-specific nutrient
application. Soil-test locations are selected before the nutrient field is
observed; conditional on acquired information, a fixed number of contiguous
management zones and a common application rate per zone are chosen to
minimize testing and expected fertilization costs subject to crop-nutrition
reliability. We first introduce a core two-stage model without explicit
homogeneity constraints. We then formulate three alternative stochastic
homogeneity specifications: a conditional expected relative-variance index,
a conditional chance constraint on the same index, and a conditional
expected within-zone variance bound. The first two preserve the relative-
variance measure of the deterministic IHDP. We give an explicit flow-based
connectivity formulation, a closed-form optimal dosing rule for fixed
partitions, and finite-scenario counterparts with probing-dependent
nonanticipativity, providing a common foundation for comparative analysis.
\end{abstract}

\section{Introduction}
\label{sec:introduction}

The delineation of site-specific management zones enables agricultural inputs to
be applied at rates that reflect spatial differences in soil properties. In the
Intra-Zone Homogeneity Delineation Problem (IHDP), a field is partitioned into a
minimum number of contiguous regions satisfying a prescribed homogeneity level,
under the assumption that the relevant soil attribute is known
\citep{VicencioMedinaIHDP}. In practice, however, nutrient availability is only
partially observed, and collecting additional soil samples entails a cost.
Furthermore, the delineation is operationally valuable because all spatial units
within a management zone can receive a common nutrient application rate.

We consider a \emph{probing-enhanced} extension of the IHDP in which soil
information is acquired endogenously before constructing the zones. The crop and
the number of management zones are specified in advance. First, the
decision-maker selects soil-testing locations. After the corresponding test
outcomes are observed, the field is partitioned into contiguous zones and a
nutrient application rate is selected for each zone. The objective is to
minimize the sum of testing costs and the expected cost of nutrient application,
while controlling within-zone nutrient heterogeneity and ensuring a prescribed
probability of meeting crop-nutrition requirements. This setting couples spatial
partitioning, information acquisition, and resource-allocation decisions in a
two-stage stochastic optimization problem.

\section{Problem definition and assumptions}
\label{sec:problem}

Let $F$ be an agricultural field partitioned into $n$ nonoverlapping geographic
units $\J=\{1,\ldots,n\}$. Each unit $j\in\J$ has area $a_j>0$; taking
$a_j=1$ represents a one-hectare grid. Spatial adjacency is represented by an
undirected graph $G=(\J,\Eset)$, where $\{j,k\}\in\Eset$ if and only if units
$j$ and $k$ share a boundary segment. We define the directed adjacency set
$\A=\{(j,k),(k,j):\{j,k\}\in\Eset\}$. The field must be partitioned into
exactly $K$ nonempty management zones, indexed by
$\I=\{1,\ldots,K\}$, with $1\leq K<n$. Each zone must induce a connected
subgraph of $G$.

A single nutrient (e.g., crop-available nitrogen) and a predetermined crop are
considered. The initial nutrient availability in unit $j$ is an uncertain
quantity $\xi_j$, expressed in nutrient-equivalent units per hectare. A known
threshold $\tau_j$ specifies the nutrient availability required in unit $j$ to
attain the target crop performance. When the crop and conditions are uniform,
one may take $\tau_j=\tau$ for all $j$. This threshold-based representation
abstracts from an explicit yield-response function: nutrient additions above
$\tau_j$ are assumed not to provide further benefits relevant to the model.
The same effective nutrient application rate is applied to all units of a
management zone. Thus, the available nutrient in unit $j$ after treatment is
$\xi_j+d_j$, where $d_j$ is the selected rate. An extension with
location-specific uptake efficiencies would require modifying this balance.

\subsection{Uncertainty, probing, and decision timing}
\label{subsec:information}

Let $\boldsymbol{\xi}=(\xi_j)_{j\in\J}$ be the vector of initial nutrient
availabilities, and let $\boldsymbol{\eta}=(\eta_j)_{j\in\J}$ be a vector of
\emph{potential} testing outcomes. The joint distribution
$P_{\boldsymbol{\xi},\boldsymbol{\eta}}$ is assumed known, including spatial
correlations between geographic units. A test may reveal nutrient availability
exactly ($\eta_j=\xi_j$) or imperfectly (for instance,
$\eta_j=\xi_j+\varepsilon_j^{\mathrm{meas}}$). The joint law is assumed to
admit regular conditional distributions, and all conditional moments used below
are finite.

In the \emph{first stage}, binary decisions $z_j$ determine whether unit $j$
is tested. The selected testing set is
\begin{equation}
  S(z)=\{j\in\J:z_j=1\}.
  \label{eq:selected-set}
\end{equation}
After the tests, the decision-maker observes only
$\boldsymbol{\eta}_{S(z)}$. In the \emph{second stage}, the realized
observation $h$ is used to determine the zone partition and nutrient rates.
The true vector $\boldsymbol{\xi}$ is not assumed to become fully observable
before these decisions; therefore, the second-stage decisions must be
measurable with respect to the acquired information
$\sigma(\boldsymbol{\eta}_{S(z)})$. There is no further corrective decision
after the actual nutrient status is realized.

The model imposes two operational requirements. First, the expected spatial
variation \emph{within each zone}, conditional on the testing information,
must remain below a prescribed tolerance. Second, each geographic unit must
meet its crop-nutrition threshold with a specified conditional probability.
These requirements are distinct: uniform nutrient addition within a zone does
not, by itself, reduce the pre-existing variation among the units of that zone.

\section{Notation and conditional coefficients}
\label{sec:notation}

Table~\ref{tab:notation} summarizes the sets, parameters, and information
variables. The second-stage optimization variables are introduced immediately
before the formulation.

\begin{table}[htbp]
\centering
\caption{Sets, parameters, and random quantities.}
\label{tab:notation}
\small
\renewcommand{\arraystretch}{1.16}
\begin{tabularx}{\linewidth}{@{}lX@{}}
\toprule
Symbol & Definition \\
\midrule
$F$, $\J$, $n$ & Field, set of geographic units, and number of units. \\
$G=(\J,\Eset)$ & Undirected geographic adjacency graph. \\
$\A$ & Directed arcs induced by $\Eset$. \\
$\I=\{1,\ldots,K\}$ & Set of management zones; $K$ is fixed. \\
$a_j>0$ & Area of unit $j$ (ha). \\
$c_j^P\geq 0$ & Cost of testing unit $j$. \\
$c^N>0$ & Cost per unit of applied nutrient. \\
$\tau_j$ & Crop-nutrition target in unit $j$ (nutrient/ha). \\
$\delta\geq0$ & Maximum permitted conditional root-mean-square intra-zone dispersion (nutrient/ha). \\
$\epsilon\in(0,1)$ & Maximum allowed conditional probability of a nutrition deficit in each unit. \\
$\bar u>0$ & Upper bound on the nutrient rate (nutrient/ha). \\
$\boldsymbol{\xi}$ & Random vector of actual initial nutrient availabilities. \\
$\boldsymbol{\eta}$ & Random vector of potential testing outcomes. \\
$z_j\in\{0,1\}$ & First-stage decision: whether unit $j$ is tested. \\
$S(z)$ & Set of units selected for testing. \\
$h$ & Realized outcome of the selected tests. \\
\bottomrule
\end{tabularx}
\end{table}

All conditional coefficients below are understood to depend on a fixed
first-stage selection $S\subseteq\J$ and a realized observation
$h$ in the support of $\boldsymbol{\eta}_S$.

\subsection{Conditional within-zone homogeneity}
\label{subsec:homogeneity}

Denote the conditional mean and covariance of the nutrient field by
\begin{equation}
\begin{split}
  \mu_j(S,h) &= \EE[\xi_j\mid\boldsymbol{\eta}_S=h],\\
  \Sigma_{jk}(S,h) &= \Cov(\xi_j,\xi_k\mid\boldsymbol{\eta}_S=h).
\end{split}
\label{eq:posterior-moments}
\end{equation}
For two geographic units, define the posterior expected squared difference
\begin{equation}
\begin{aligned}
 D_{jk}(S,h)
 &:=\EE\!\left[(\xi_j-\xi_k)^2\mid\boldsymbol{\eta}_S=h\right]\\
 &=\bigl(\mu_j(S,h)-\mu_k(S,h)\bigr)^2
   +\Sigma_{jj}(S,h)+\Sigma_{kk}(S,h)-2\Sigma_{jk}(S,h).
\end{aligned}
\label{eq:pairwise-difference}
\end{equation}
These coefficients account for both the estimated spatial pattern and the
remaining uncertainty, including spatial dependence. In particular,
$D_{jj}(S,h)=0$.

For a nonempty zone $i$, let
\begin{equation}
  \rho_i(x)=\sum_{j\in\J}x_{ji},\qquad
  \bar\xi_i(x,\boldsymbol{\xi})
   =\frac{\sum_{j\in\J}x_{ji}\xi_j}{\rho_i(x)}.
  \label{eq:zone-mean}
\end{equation}
We measure its realized intra-zone dispersion by the (population-normalized)
within-zone variance. Geographic units are weighted equally in this statistic,
as appropriate for the one-hectare grid; an area-weighted variant may be used
for unequal spatial cells. The variance is
\begin{equation}
  V_i(x,\boldsymbol{\xi})
  =\frac{1}{\rho_i(x)}\sum_{j\in\J}
       x_{ji}\bigl(\xi_j-\bar\xi_i(x,\boldsymbol{\xi})\bigr)^2.
  \label{eq:within-zone-var}
\end{equation}
The identity relating variance to pairwise squared differences yields
\begin{equation}
  \EE\!\left[V_i(x,\boldsymbol{\xi})
       \mid\boldsymbol{\eta}_S=h\right]
  =\frac{1}{2\rho_i(x)^2}
       \sum_{j\in\J}\sum_{k\in\J}
          D_{jk}(S,h)x_{ji}x_{ki}.
  \label{eq:conditional-zone-var}
\end{equation}
We require this quantity to be at most $\delta^2$ for every zone. This
zone-by-zone requirement deliberately differs from the \emph{global} relative
variance (RV) measure used in the original IHDP. It is chosen because the
same nutrient rate will be applied throughout each delineated zone.

\subsection{Conditional nutrient sufficiency}
\label{subsec:quantile}

For each unit $j$, define the conditional lower-tail threshold
\begin{equation}
  q_{\epsilon,j}(S,h)
  :=\sup\Bigl\{t\in\mathbb R:
    \PP\bigl(\xi_j<t\mid\boldsymbol{\eta}_S=h\bigr)
       \leq\epsilon\Bigr\}.
  \label{eq:quantile-definition}
\end{equation}
We assume this threshold is finite. The strict inequality inside the
probability makes the definition appropriate for both continuous and discrete
posterior distributions. When the conditional distribution is continuous and
strictly increasing at the relevant level, it is the usual lower
$\epsilon$-quantile. For a fixed application rate $d_j$, the following
conditions are equivalent:
\begin{equation}
\begin{aligned}
 \PP\bigl(\xi_j+d_j\geq\tau_j
      \mid\boldsymbol{\eta}_S=h\bigr)&\geq1-\epsilon
 \\[-1pt]
 &\Longleftrightarrow\quad
 d_j\geq\tau_j-q_{\epsilon,j}(S,h).
\end{aligned}
\label{eq:chance-equivalence}
\end{equation}
Thus, conditional nutritional reliability can be enforced through a
deterministic lower bound on $d_j$ once the posterior is known. These are
\emph{individual} chance constraints: a probability $1-\epsilon$ is required
for each unit separately, not jointly for the entire field.

\section{Two-stage stochastic optimization formulation}
\label{sec:formulation}

\subsection{Decisions and objective}
\label{subsec:decisions}

For a given testing set $S$ and observation $h$, the second-stage variables
are
\begin{itemize}
\item $x_{ji}\in\{0,1\}$: one if unit $j$ belongs to management zone $i$;
\item $r_{ji}\in\{0,1\}$: one if unit $j$ is selected as the root of zone $i$;
\item $u_i\in[0,\bar u]$: nutrient application rate shared by all units in zone $i$;
\item $d_j\in[0,\bar u]$: nutrient application rate received by unit $j$;
\item $g_{ij}\geq0$: artificial flow sent from a super-source to unit $j$ for zone $i$;
\item $f_{ijk}\geq0$: artificial flow on directed arc $(j,k)\in\A$ for zone $i$.
\end{itemize}
The flow and root variables serve only to certify spatial connectivity;
they do not represent physical movement of fertilizer.

The first-stage objective is
\begin{equation}
\boxed{\displaystyle
\min_{z\in\{0,1\}^{n}}
 \left\{
   \sum_{j\in\J} c_j^P z_j
    +\EE_{\boldsymbol{\eta}_{S(z)}}\!\left[
       Q\bigl(S(z),\boldsymbol{\eta}_{S(z)}\bigr)
     \right]
 \right\}.}
\label{eq:first-stage}
\end{equation}
The first term is the cost of acquiring information. The recourse function
$Q(S,h)$ is the minimum nutrient-application cost when the available
information is $\boldsymbol{\eta}_{S}=h$. In particular,
$S=\varnothing$ is permitted and corresponds to no additional soil testing.

\subsection{Second-stage model}
\label{subsec:recourse}

For fixed $(S,h)$, the second-stage optimization problem is
\allowdisplaybreaks
\begin{align}
Q(S,h)=\min_{x,r,g,f,u,d}\quad
 &c^N\sum_{j\in\J}a_jd_j
 &&\label{eq:recourse-objective}\\
\text{s.t.}\quad
 &\sum_{i\in\I}x_{ji}=1
 &&\forall j\in\J,\label{eq:assign}\\
 &\sum_{j\in\J}x_{ji}\geq1
 &&\forall i\in\I,\label{eq:nonempty}\\
 &\sum_{j\in\J}r_{ji}=1
 &&\forall i\in\I,\label{eq:one-root}\\
 &r_{ji}\leq x_{ji}
 &&\forall i\in\I,\ j\in\J,\label{eq:root-located}\\
 &\sum_{j\in\J}g_{ij}=\sum_{j\in\J}x_{ji}
 &&\forall i\in\I,\label{eq:source-balance}\\
 &0\leq g_{ij}\leq n\,r_{ji}
 &&\forall i\in\I,\ j\in\J,\label{eq:source-capacity}\\
 &\sum_{k:(k,j)\in\A}f_{ikj}+g_{ij}
       -\sum_{k:(j,k)\in\A}f_{ijk}=x_{ji}
 &&\forall i\in\I,\ j\in\J,\label{eq:flow-balance}\\
 &0\leq f_{ijk}\leq (n-1)x_{ji}
 &&\forall i\in\I,\ (j,k)\in\A,\label{eq:flow-tail}\\
 &f_{ijk}\leq(n-1)x_{ki}
 &&\forall i\in\I,\ (j,k)\in\A,\label{eq:flow-head}\\
 &d_j-u_i\leq\bar u(1-x_{ji})
 &&\forall i\in\I,\ j\in\J,\label{eq:dose-link-1}\\
 &u_i-d_j\leq\bar u(1-x_{ji})
 &&\forall i\in\I,\ j\in\J,\label{eq:dose-link-2}\\
 &\sum_{j,k\in\J}D_{jk}(S,h)x_{ji}x_{ki}
     \leq2\delta^2\!\left(\sum_{j\in\J}x_{ji}\right)^{\!2}
 &&\forall i\in\I,\label{eq:homogeneity}\\
 &d_j\geq\tau_j-q_{\epsilon,j}(S,h)
 &&\forall j\in\J,\label{eq:nutrition}\\
 &x_{ji},r_{ji}\in\{0,1\}
 &&\forall i\in\I,\ j\in\J,\label{eq:binary-domain}\\
 &f_{ijk}\geq0,\quad g_{ij}\geq0
 &&\forall i\in\I,\ (j,k)\in\A,\ j\in\J,\label{eq:flow-domain}\\
 &0\leq u_i\leq\bar u,\quad0\leq d_j\leq\bar u
 &&\forall i\in\I,\ j\in\J.\label{eq:dose-domain}
\end{align}
We set $Q(S,h)=+\infty$ if no feasible second-stage decision exists.
The main formulation therefore selects only testing designs that admit a
feasible recourse decision for almost every observation. This is an explicit
feasibility assumption of the base model: a fixed $K$ and a stringent
homogeneity tolerance may otherwise prevent a feasible partition for some
observations. In applications without relatively complete recourse, one can
introduce explicitly penalized violations or revise the admissible testing
designs. The bound $\bar u$ must also be chosen consistently with the
nutritional targets and the support of the posterior distributions.

\subsection{Interpretation of the constraints}
\label{subsec:interpretation}

The objective \eqref{eq:first-stage} minimizes the \emph{expected total cost}
of probing and nutrient application. The recourse objective
\eqref{eq:recourse-objective} charges for the amount applied to each geographic
unit, taking its area into account. Hence, when all $a_j=1$, the cost is
proportional to the sum of applied nutrient rates over all hectares.

Constraint \eqref{eq:assign} assigns every unit to exactly one zone, while
\eqref{eq:nonempty} enforces exactly $K$ nonempty zones, as their number is
fixed in advance. Constraints \eqref{eq:one-root}--\eqref{eq:root-located}
select one geographic root per zone. Constraints
\eqref{eq:source-balance}--\eqref{eq:flow-head} enforce \emph{connectivity}:
the artificial super-source supplies one unit of flow to every selected
geographic unit, and positive flow may traverse an adjacency arc only when
both endpoints belong to the same zone. Since only the root can receive flow
from the super-source, every unit assigned to a zone must be connected to
its root through units of that same zone. Consequently, each zone induces a
connected subgraph of $G$.

Constraints \eqref{eq:dose-link-1}--\eqref{eq:dose-link-2} impose the
\emph{uniform nutrient rate} requirement. If $x_{ji}=1$, they force
$d_j=u_i$. If $x_{ji}=0$, they are inactive because both rates belong to
$[0,\bar u]$. Thus, $d_j$ is precisely the common rate of the zone to
which unit $j$ is assigned; no bilinear product $x_{ji}u_i$ is needed in the
objective.

Constraint \eqref{eq:homogeneity} is the within-zone homogeneity
requirement. Because each zone is nonempty, Equation
\eqref{eq:conditional-zone-var} shows that this quadratic constraint is
exactly equivalent to
$\EE[V_i(x,\boldsymbol{\xi})\mid\boldsymbol{\eta}_S=h]\leq\delta^2$.
The tolerance $\delta$ is specified by the decision-maker and applies to
each zone separately. Constraint \eqref{eq:nutrition} enforces the
conditional chance constraint
$\PP(\xi_j+d_j\geq\tau_j\mid\boldsymbol{\eta}_S=h)\geq1-\epsilon$
for each geographic unit. Finally,
\eqref{eq:binary-domain}--\eqref{eq:dose-domain} define the variable
domains. The homogeneity restrictions are generally nonconvex in the binary
assignment decisions, so the proposed recourse problem is not claimed to be
a mixed-integer linear program.

\Needspace{16\baselineskip}
\subsection{Optimal nutrient rate for a fixed partition}
\label{subsec:optimal-dose}

The model has an elementary but useful structure: once the partition has
been fixed, the optimal nutrient rates can be obtained analytically.

\begin{proposition}\label{prop:optimal-dose}
Fix $S$, $h$, and a partition $\Z=\{Z_1,\ldots,Z_K\}$ satisfying the
assignment, nonemptiness, connectivity, and conditional homogeneity
requirements. Suppose the nutrient cost $c^N$ and all areas $a_j$ are
positive, and the resulting nutrient rates do not exceed $\bar u$. Then
the optimal rate for zone $i$ is
\begin{equation}
  u_i^*(S,h,\Z)
  =\max\left\{0,\,\max_{j\in Z_i}
       \bigl(\tau_j-q_{\epsilon,j}(S,h)\bigr)\right\}.
  \label{eq:closed-form-dose}
\end{equation}
\end{proposition}

\begin{proof}
For $j\in Z_i$, constraints
\eqref{eq:dose-link-1}--\eqref{eq:dose-link-2} imply $d_j=u_i$.
Constraint \eqref{eq:nutrition} consequently requires
$u_i\geq\tau_j-q_{\epsilon,j}(S,h)$ for every $j\in Z_i$, as well as
$u_i\geq0$. The nutrient cost of zone $i$ equals
$c^N u_i\sum_{j\in Z_i}a_j$, which is strictly increasing in $u_i$.
The smallest feasible rate is therefore optimal.
\end{proof}

Accordingly, the operational cost of a fixed partition is
\begin{equation}
  C^N(\Z\mid S,h)
  =c^N\sum_{i=1}^{K}\left(\sum_{j\in Z_i}a_j\right)
       \max\left\{0,\max_{j\in Z_i}
           \bigl(\tau_j-q_{\epsilon,j}(S,h)\bigr)\right\}.
  \label{eq:partition-cost}
\end{equation}
This expression makes the economic role of delineation explicit: a single
low-nutrient unit can determine the rate applied to an entire management
zone. Grouping units with similar nutrient needs can therefore reduce
fertilizer expenditures even when the number of zones is fixed.

\section{Finite-scenario representation and endogenous nonanticipativity}
\label{sec:scenario}

For computation, suppose the joint distribution is represented by a finite
scenario set $\Om$. Each scenario $\omega\in\Om$ has probability
$\pi_\omega>0$ with $\sum_{\omega\in\Om}\pi_\omega=1$, and specifies
$(\boldsymbol{\xi}^{\omega},\boldsymbol{\eta}^{\omega})$. Scenario-dependent
second-stage variables are denoted by
$(x^{\omega},r^{\omega},g^{\omega},f^{\omega},u^{\omega},d^{\omega})$.
The first-stage testing decisions $z$ are common to all scenarios.

For any two scenarios, let
\begin{equation}
  \Delta_{\omega\nu}
   =\{j\in\J:\eta_j^{\omega}\neq\eta_j^{\nu}\},\qquad
  m_{\omega\nu}(z)=\sum_{j\in\Delta_{\omega\nu}}z_j.
  \label{eq:distinguishability}
\end{equation}
Thus, $m_{\omega\nu}(z)=0$ if and only if the selected tests cannot
distinguish scenarios $\omega$ and $\nu$. The corresponding information
class and its probability are
\begin{equation}
\begin{aligned}
 \mathcal C_\omega(z)
   &=\{\nu\in\Om:m_{\omega\nu}(z)=0\},\\
 P_\omega(z)&=\sum_{\nu\in\mathcal C_\omega(z)}\pi_\nu.
\end{aligned}
\label{eq:information-class}
\end{equation}
For scenario $\omega$, conditional coefficients are computed from
this information class. In particular,
\begin{equation}
  D_{jk}^{\omega}(z)
   =\frac{1}{P_\omega(z)}
       \sum_{\nu\in\mathcal C_\omega(z)}
          \pi_\nu(\xi_j^{\nu}-\xi_k^{\nu})^2,
  \label{eq:scenario-dispersion}
\end{equation}
while $q_{\epsilon,j}^{\omega}(z)$ is the lower-tail threshold
\eqref{eq:quantile-definition} of the discrete distribution placing mass
$\pi_\nu/P_\omega(z)$ on $\xi_j^{\nu}$ for
$\nu\in\mathcal C_\omega(z)$.

The finite-scenario counterpart minimizes
\begin{equation}
 \min\quad
   \sum_{j\in\J}c_j^Pz_j
   +c^N\sum_{\omega\in\Om}\pi_\omega
      \sum_{j\in\J}a_jd_j^{\omega},
 \label{eq:scenario-objective}
\end{equation}
subject to a copy of the second-stage constraints
\eqref{eq:assign}--\eqref{eq:dose-domain} for each $\omega\in\Om$,
using $D_{jk}^{\omega}(z)$ and $q_{\epsilon,j}^{\omega}(z)$ in
\eqref{eq:homogeneity} and \eqref{eq:nutrition}, together with the
following \emph{probing-dependent nonanticipativity} conditions:
\begin{equation}
\begin{aligned}
 \bigl|x_{ji}^{\omega}-x_{ji}^{\nu}\bigr|
   &\leq m_{\omega\nu}(z)
   &&\forall \omega,\nu\in\Om,\ j\in\J,\ i\in\I,\\
 \bigl|u_i^{\omega}-u_i^{\nu}\bigr|
   &\leq \bar u\,m_{\omega\nu}(z)
   &&\forall \omega,\nu\in\Om,\ i\in\I.
\end{aligned}
\label{eq:NA}
\end{equation}

When $m_{\omega\nu}(z)=0$, the scenarios yield identical observations,
and \eqref{eq:NA} requires identical partitions and zone
application rates. When at least one selected test
reveals different outcomes, these constraints permit different decisions.
Each absolute-value inequality in \eqref{eq:NA} is equivalent to
two linear inequalities. The remaining variables $r,g,f$ provide mathematical connectivity
certificates and need not be made nonanticipative: identical zone
assignments in indistinguishable scenarios admit a common connectivity
certificate. Likewise, the physical unit-level dose $d_j$ is identical
whenever the associated $x$ and $u$ decisions coincide.

Because the posterior coefficients in
\eqref{eq:scenario-dispersion} and the conditional quantiles depend on
$z$ through $\mathcal C_\omega(z)$, this is an exact mathematical
representation of a \emph{finite-distribution} probing-enhanced problem,
but is not directly a mixed-integer linear formulation. Deriving tractable
reformulations, decomposition algorithms, or tailored branching rules is a
separate methodological step. Furthermore, if the original signals are
continuous, care is needed when constructing the finite scenario set:
independently sampled continuous outcomes are almost surely distinct,
and treating each distinct sample as a separate observation would
artificially provide perfect identification of the sampled scenario.
Signal discretization or a conditional-distribution approximation is
therefore required for a meaningful scenario-based implementation.

\section{Modeling scope and possible extensions}
\label{sec:extensions}

The proposed model isolates the value of information for nutrient management
under a fixed crop and a fixed number of connected zones. Its homogeneity
criterion applies to the \emph{initial} nutrient field. A uniform additive
rate cannot alter within-zone spatial variation under the assumed nutrient
balance, but an appropriate delineation avoids combining units requiring
substantially different nutrient rates. The chance constraints guarantee a
minimum nutritional condition; they do not explicitly maximize expected
crop yield. If measured yield-response curves are available, one may
replace the nutritional thresholds by a yield-based objective or yield
reliability constraints. Other extensions include multiple nutrients,
zone-specific application setup costs, physical constraints on fertilizer
machinery, adaptive sequential testing, and explicit penalties for
unattainable nutrition or homogeneity targets.


\bibliographystyle{plainnat}
\bibliography{references}

\end{document}
